

\subsection{Adversarial Injection}
\label{sec:rq3}
This intervention assesses for attentive listening (Q3).

\noindent \textbf{Setup:} We generate an \textit{Adversarial Dataset} featuring conflicting linguistic and acoustic cues. We have two setups. In the first setup, we inject a speech signal (generated by CosyVoice TTS~\cite{du2024cosyvoice}) that gives utters a wrong answer multiple times. In the second setup we inject a speech signal that utters the correct answer multiple times. In both setups the power of the adversarial speech is comparable or lower than the input audio. We intend to measure whether the model's rationale remains grounded in the audio (which is required to answer the question) or is misled by the text transcription.

\noindent \textbf{Results:} In Figure \ref{fig:adv_results}, we observe that the adversarial injections in fact result in very significant performance reductions when a wrong answer is injected in the audio. We observe for instance, that for the Sakura-Animal dataset, the performance drops from $\mathord{\sim}75\%$ to $\mathord{\sim}25\%$ for AF3. QWEN seems to be slightly more robust, but it still suffers from significant performance alterations because of the adversarial interventions. These results suggest that LALMs may rely more heavily on linguistic cues than on the underlying acoustic signal.

